\SourceLeaf{21}
\par
Man sieht leicht, dass so die gesamte Oberfläche des Sphäroids erfasst wird, wenn $p$ von $0$ bis $180^0$ und $q$ von $0$ bis $360^0$ läuft. Ferner erhalten wir
\[
\begin{aligned}
\lambda&=-A\sin p, &\quad\lambda'&=0\\
\mu&=B\cos p.\cos q, &\quad\mu'&=-B\sin p.\sin q\\
\nu&=C\cos p.\sin q, &\quad\nu'&=C\sin p.\cos q
\end{aligned}
\]
\[
\begin{aligned}
\mu\nu'-\nu\mu'&=BC\cos p.\sin p=ABC\sin p.\frac{x}{AA}\\
\nu\lambda'-\lambda\nu'&=AC\sin p^2.\cos q=ABC\sin p.\frac{y}{BB}\\
\lambda\mu'-\mu\lambda'&=AB\sin p^2.\sin q=ABC\sin p.\frac{z}{CC}
\end{aligned}
\]
\par
Da $\sin p$ innerhalb der hier betrachteten Grenzen überall positiv ist, ist also zu setzen
\[
\mathrm{d}s=\mathrm{d}p.\mathrm{d}q.ABC.\psi.\sin p
\]
\par
Wenden wir diese Formeln auf den zweiten Lehrsatz an, so ergibt sich das Volumen des Körpers oder (wenn wir die Dichte $=1$ setzen) seine Masse
\[
=\int\!\int\mathrm{d}p.\mathrm{d}q.ABC.\cos p^2.\sin p
\]
oder, wenn man zuerst nach $q$ integriert,
\[
=2\pi\int\mathrm{d}p.ABC.\cos p^2.\sin p=\tfrac12\pi ABC\int\mathrm{d}p.(\sin p+\sin3p)
\]
Dieses Integral ist von $p=0$ bis $p=180^0$ zu erstrecken. Daraus ergibt sich $\tfrac43\pi ABC$, wie auch anderweitig bekannt ist.

\Article{13}
\par
Um die Anziehung zu bestimmen, die das Sphäroid auf einen beliebigen Punkt ausübt, nehmen wir an, dass die Anziehung jedes Elements dem Quadrat seines Abstands vom angezogenen Punkt umgekehrt proportional ist. Dann gilt $fr=\frac{1}{rr}$, $Fr=-\frac{1}{r}$, $\Phi r=r$. Die Anziehung des gesamten Sphäroids in der zur $x$-Koordinatenachse parallelen und ihr entgegengesetzten Richtung sei $=X$, und wir setzen $X=ABC\xi$. Nach dem dritten Lehrsatz gilt also
\[
X=\int\!\int\mathrm{d}p.\mathrm{d}q\frac{BCx\sin p}{rA}=\int\!\int\mathrm{d}p.\mathrm{d}q\frac{BC\cos p.\sin p}{r}
\]
und folglich
